\hypertarget{ej11}{}
\noindent \textbf{Ejercicio 11.} Sean $P(x:\mathbb{Z})$ y $Q(x:\mathbb{Z}, y:\mathbb{Z})$ dos predicados cualesquiera que nunca se indefinen. Escribir el predicado asociado a cada uno de los siguientes enunciados:

\begin{enumerate}[label=\alph*)]
    \item “Si un entero cumple P, entonces existe un entero distinto tal que juntos cumplen Q“
    \item “Existe un par de enteros tal que cumplen Q y ninguno de ambos cumple P“
    \item “Si un par de enteros cumplen Q, entonces al menos uno de ellos cumple P“
    \item “Si un entero cumple P, no existe ningún entero tal que juntos cumplan Q“
\end{enumerate}

\vspace{0.2cm}
{\color{gray}\hrule}
\vspace{0.4cm}

\textbf{Resolución:}

\begin{enumerate}[label=\textbf{\alph*)}]

    \item \textbf{“Si un entero cumple P, entonces existe un entero distinto tal que juntos cumplen Q“} \\
    La frase "Si un entero..." establece una condición general, por lo que usamos un cuantificador universal ($\forall$). Para decir que es "distinto", exigimos explícitamente que la nueva variable sea diferente a la primera ($x \neq y$).
    \[ (\forall x: \mathbb{Z})(P(x) \rightarrow (\exists y: \mathbb{Z})(x \neq y \wedge Q(x, y))) \]

    \vspace{0.4cm}
    \item \textbf{“Existe un par de enteros tal que cumplen Q y ninguno de ambos cumple P“} \\
    Acá buscamos la existencia de dos variables. La palabra "par" sugiere fuertemente que son dos elementos distintos, por lo que es una buena práctica (y más riguroso) agregar $x \neq y$. Que "ninguno" cumpla $P$ significa que tanto $x$ como $y$ hacen falso al predicado $P$.
    \[ (\exists x, y: \mathbb{Z})(x \neq y \wedge Q(x, y) \wedge \neg P(x) \wedge \neg P(y)) \]

    \vspace{0.4cm}
    \item \textbf{“Si un par de enteros cumplen Q, entonces al menos uno de ellos cumple P“} \\
    Nuevamente, el "Si un par..." actúa como una regla general, por lo que usamos un $\forall$. Si la condición (cumplir $Q$) se da, entonces $P(x)$ o $P(y)$ debe ser verdadero (disyunción).
    \[ (\forall x, y: \mathbb{Z})(x \neq y \wedge Q(x, y) \rightarrow (P(x) \vee P(y))) \]
    \textit{Nota: Dependiendo de cuán estricto sea el docente con la palabra "par", el $x \neq y$ en el antecedente puede omitirse, pero ponerlo demuestra que entendemos que hablamos de dos enteros diferentes.}

    \vspace{0.4cm}
    \item \textbf{“Si un entero cumple P, no existe ningún entero tal que juntos cumplan Q“} \\
    Arrancamos con el $\forall$ general para el primer entero. Luego, en el consecuente, negamos literalmente la existencia de un segundo entero que cumpla la relación $Q$ con el primero.
    \[ (\forall x: \mathbb{Z})(P(x) \rightarrow \neg (\exists y: \mathbb{Z}) Q(x, y)) \]
    \textit{Alternativa equivalente (pasando la negación hacia adentro):} 
    \[ (\forall x: \mathbb{Z})(P(x) \rightarrow (\forall y: \mathbb{Z}) \neg Q(x, y)) \]

\end{enumerate}

\vspace{0.5cm}
\hrule
\vspace{0.5cm}